\documentclass[11pt]{letter}

\usepackage[margin=1in]{geometry}
\usepackage{amsmath}

\signature{Stephen Nelson\\Field Dynamics Research}

\begin{document}

\begin{letter}{%
The Editor\\
\textit{Classical and Quantum Gravity}%
}

\opening{Dear Editor,}

I am submitting the enclosed manuscript, \textit{The Other Half of
Gravity: Torsion and the Two-Chamber Structure of the Tetrad Frame},
for consideration as a Letter in \textit{Classical and Quantum Gravity}.
The paper is a formal companion to \textit{The Coordinate-Space Geometry
of the Tetrad Frame} (CQG-115428), currently under review at your journal.

That companion letter proved that the tetrad equation produces two
bounded tetrahedral regions related by spatial parity: TetA (the
identity component) and TetB (the parity component). The present letter
asks what becomes of this two-chamber structure at the level of the
affine connection.

The central result is a lemma: a torsion-free connection cannot
distinguish TetA from TetB. The affine connection carries $K^3 = 64$
components; the 40 Christoffel symbols are built entirely from the
metric, which is invariant under the parity exchanging the two chambers,
while the 24 torsion components depend on the frame field directly and
are therefore the necessary carrier of the two-chamber distinction. The
zero-torsion condition collapses that distinction. A second result links
the antisymmetric count~$24 = D \times 2K$ to a coupling polynomial
that organises the standard count structure of the tetrad formalism,
and this identity holds exclusively at~$K = 4$.

A supplementary note derives the TetB field fraction $1/(x{+}1)$
from the exchange symmetry of the two chambers. The ratio of removed
to retained field response is $1/x$, which is algebraically identical
to the dark-to-baryon mass ratio of the cold dark matter framework,
following from four lines of algebra with no halo model and no free
parameters.

The manuscript does not propose a modification of general relativity.
All claims are either proved locally, derived from the companion
geometry, or cited from standard references (Misner, Thorne and Wheeler;
Hehl \textit{et al.}\ 1976; Cartan 1922). A validation script confirms
52/52 symbolic claims.

Because the present letter builds directly on the geometry established
in CQG-115428, I would be grateful if the same referee could be
considered for both manuscripts, as this would avoid the need to
re-establish the companion result within the present review.

The manuscript has not been submitted elsewhere and is not under
consideration at any other journal.

\closing{Yours sincerely,}

\end{letter}

\end{document}
